\begin{frame}{Estrategias de Reforma Tributaria}
\begin{itemize}
%    \item \textbf{Meta}: aumentar recaudo permanente en al menos 2\% del PIB.
    \item \textbf{Beneficios Tributarios}:
    \begin{itemize}
        %\item Revisar y desmontar beneficios tributarios más costosos (1,3\% PIB).
        %\item Elevar tasa mínima efectiva al 20\% sobre base tributaria.
        \item Reformar Estatuto Tributario para explicitar beneficios y eliminarlos por defecto.
        \item Eventual eliminación de beneficios cuyo gasto tributario no esté cuantificado dentro del PGN.
    \end{itemize}
\end{itemize}
\end{frame}

\begin{frame}{Estrategias de Reforma Tributaria}
\begin{itemize}
    \item \textbf{Renta personas jurídicas}:
    \begin{itemize}
        \item Revisar y desmontar beneficios tributarios más costosos (Hasta 1,3\% PIB).
        \item Elevar tasa mínima efectiva al 20\% sobre base tributaria.
        \item Reducir tarifa general mediante eliminación de rentas de destinación específica y traslado de competencias.
        \item Combatir uso indebido de empresas para gastos personales (Hasta 1\% PIB).
    \end{itemize}
\end{itemize}
\end{frame}

\begin{frame}{Estrategias de Reforma Tributaria}
\begin{itemize}    
    \item \textbf{Renta personas naturales}:
    \begin{itemize}
        \item Evitar mayor carga sobre rentas laborales.
        \item Mayor control a rentas no laborales con deducciones inusuales.
        \item Eliminar exención a pensiones.
        \item Eliminar doble tributación de dividendos
    \end{itemize}
\end{itemize}
\end{frame}

\begin{frame}{Estrategias de Reforma Tributaria}
\begin{itemize}
    \item \textbf{IVA} ({Hasta \raisebox{0.2ex}{\texttildelow}}5\% PIB):
    \begin{itemize}
        \item Ampliar base: pasar exentos a excluidos y eliminar tarifa del 5\%.
        \item Considerar bajar tarifa general del 19\% al 16\% con base más amplia.
        \item Implementar Registro Único de Ingresos para compensaciones focalizadas.
    \end{itemize}
    \item \textbf{Otros impuestos}:
    \begin{itemize}
        \item Ampliar base del GMF a transacciones móviles (0,2\% PIB).
        \item Aumentar ImpoCarbono %y/o eliminar subsidio a combustibles (0,6\% PIB).
        \item Eliminar Impoconsumo y SIMPLE por baja eficiencia (-0,3\% PIB).
        \item Suspender Impuesto al Patrimonio.
    \end{itemize}
\end{itemize}
\end{frame}



\begin{frame}{Estrategias de Reforma Tributaria}
\begin{itemize}    
    \item \textbf{Control y equidad}:
    \begin{itemize}
        %\item Combatir uso indebido de empresas para gastos personales (1\% PIB).
        \item Modernizar tecnología aduanera para reducir contrabando (0,5\% PIB).
    \end{itemize}
    \item \textbf{Economía Política}:
    \begin{itemize}    
        \item Discusión con las cortes sobre ``derechos adquiridos'' de las personas jurídicas.
        \item Ley de Cabildeo.
        \item Acción colectiva de los gremios.
    \end{itemize}
    \item \textbf{Reforma Tributaria Territorial}
\end{itemize}
\end{frame}
